\section{The depth-three circuit lower bound}\label{sec:lower-bound}

We now combine the two structural lemmas with the hard slices.
The order of choices is useful: fix the proposed circuit exponent \(A\),
then choose constants \(B,b\), and only then let the input length grow.
All choices of parameter bits remain restrictions of the language already
defined.

\begin{proof}[Proof of Theorem~\ref{thm:main}]
Proposition~\ref{prop:polytime} supplies the single polynomial-time
machine.  Fix \(A>0\), and put
\[
\begin{gathered}
 d=\floor{n/5},\qquad s=3A,\qquad B=64(s+1),\\
 p=\frac{B}{\sqrt d},\qquad m_0=B\sqrt d,\qquad
 k=\ceil{3(s+1)\sqrt d}.
\end{gathered}
\]
All subsequent assertions hold for every sufficiently large \(n\),
with the threshold allowed to depend on \(A\).
In particular, \(0<p\le1/2\) and \(n\le9d\), so
\begin{equation}\label{eq:circuit-size-scale}
 2^{A\sqrt n}\le2^{s\sqrt d}.
\end{equation}

\paragraph{Choose a constant target width.}
Since \(k\le[3(s+1)+1]\sqrt d\) for \(d\ge1\), we have
\[
 \frac{2pk}{1-p}\le Q,\qquad
 Q=4B[3(s+1)+1].
\]
Choose an integer \(b\ge1\), depending only on \(A\), so large that
\begin{equation}\label{eq:factorial-tail-choice}
 \sum_{\ell>b}\frac{Q^\ell}{\ell!}\le\frac12.
\end{equation}
Such a choice exists: once \(\ell+1\ge2Q\), successive terms have
ratio at most \(1/2\), and the tails tend to zero.
Using \(\binom{k}{\ell}\le k^\ell/\ell!\), the parameter of
Lemma~\ref{lem:restriction} satisfies
\begin{equation}\label{eq:theta-half}
 \theta=\sum_{\ell=b+1}^k
       \binom{k}{\ell}\left(\frac{2p}{1-p}\right)^\ell
 \le\sum_{\ell>b}\frac{Q^\ell}{\ell!}\le\frac12.
\end{equation}
This is the only place where the desired circuit exponent affects the
target width.

\paragraph{Transfer correlation to width \(k\).}
Fix a setting of the non-data bits supplied by
Proposition~\ref{prop:good-slices}, and let \(g:\bits^d\to\{-1,1\}\)
be the sign of that slice.
Call a restriction good if it lies outside the failure event in
\eqref{eq:good-slices}.  On such a restriction,
\[
 \corr_b(g_{T,\rho})
 \le2^{-|T|/8}\le2^{-m_0/16}.
\]
The exceptional probability is at most \(4\cdot2^{-m_0/16}\).

For any fixed \(H\in\mathcal C_k([d])\),
Lemma~\ref{lem:restriction} gives a nonnegative \(W_H\) with
\(\E W_H\le2\) and
\(\abs{\E_z g_{T,\rho}H_{T,\rho}}\le W_H\corr_b(g_{T,\rho})\).
On an exceptional restriction the same absolute correlation is at
most one, since \(g\) is a sign and \(H\) is an indicator.
Equation~\eqref{eq:restriction-expectation} therefore yields
\begin{align}
 \abs{\E_xg(x)H(x)}
 &\le 2^{-m_0/16}\E[W_H\one_{\mathrm{good}}]
        +\Pr[\mathrm{not\ good}]\notag\\
 &\le6\cdot2^{-m_0/16}.
 \label{eq:wide-correlation}
\end{align}
The good event and \(W_H\) need not be independent.
Moreover, the setting of the non-data bits was fixed before \(H\)
was chosen: the good event controls all width-\(b\) tests simultaneously,
and the restriction lemma applies to each width-\(k\) test.
Thus \eqref{eq:wide-correlation} holds for every such \(H\)
for this one slice.

The constant true test is a width-\(k\) CNF.  Applying
\eqref{eq:wide-correlation} to it gives
\(\abs{\E g}\le6\cdot2^{-m_0/16}\).
The accepted set of the slice consequently has density
\begin{equation}\label{eq:accepted-density}
 \Pr[g=1]=\frac{1+\E g}{2}\ge\frac14
\end{equation}
for every sufficiently large \(n\).

\paragraph{Remove wide bottom clauses and count every gate.}
Suppose, for a contradiction, that the \(n\)-bit membership function
has an OR--AND--OR circuit with \(S\le2^{A\sqrt n}\) total gates.
Delete every gate with no directed path to the output, so every
remaining middle gate feeds the top OR.
Fix its non-data inputs to the setting just chosen.
This does not increase its number of gates.
There are at most \(S\) middle AND gates.  Each is the indicator
of a CNF \(H\) whose clauses are the outputs of its distinct bottom
OR gates, and hence it has at most \(S\) clauses.
A bottom gate shared between several middle gates still supplies at
most one clause to each; repeated wires can be removed.

Normalize the clauses after the inputs are fixed.  A middle gate with a
false clause computes zero and can be omitted.  For any remaining
middle CNF \(H\), delete every clause of width greater than \(k\),
and denote the resulting width-\(k\) CNF by \(H'\).
Then \(H'\ge H\).  A normalized clause of width \(\ell>k\)
is false on a fraction \(2^{-\ell}\le2^{-k}\) of the cube,
so a union bound over its at most \(S\) deleted clauses gives
\begin{equation}\label{eq:clause-truncation}
 \E(H'-H)\le S2^{-k}.
\end{equation}
Every input accepted by \(H\) is accepted by the top OR, which computes
the slice exactly.  Hence \(gH=H\) pointwise, and
\begin{align}
 \E H=\E(gH)
 &\le \abs{\E(gH')}+\E(H'-H)\notag\\
 &\le 6\cdot2^{-m_0/16}+S2^{-k}.
 \label{eq:middle-gate-density}
\end{align}
Here \(\abs{g}=1\) bounds the truncation error, and
\eqref{eq:wide-correlation} bounds the first term.

The accepted set is the union of the sets accepted by the middle gates.
Summing \eqref{eq:middle-gate-density} over at most \(S\) of them,
and using \eqref{eq:circuit-size-scale}, gives
\begin{align}
 \Pr[g=1]
 &\le6S2^{-m_0/16}+S^2 2^{-k}\notag\\
 &\le6\cdot2^{-(3s+4)\sqrt d}
           +2^{-(s+3)\sqrt d}\longrightarrow0.
 \label{eq:final-contradiction}
\end{align}
Indeed, \(m_0/16=4(s+1)\sqrt d\) and
\(k\ge3(s+1)\sqrt d\).
This contradicts \eqref{eq:accepted-density}.
For the fixed \(A\), every sufficiently large \(n\) therefore satisfies
\(S_3(f_n)>2^{A\sqrt n}\).
Since \(A>0\) was arbitrary and the language and its machine are fixed,
Theorem~\ref{thm:main} follows.
\end{proof}
